\documentclass[12pt]{article}
\usepackage[utf8]{inputenc}
\usepackage[T1]{fontenc}
\usepackage{hyperref}
\title{North Awyu Language Resources for K-12 Curriculum Development}
\author{Compiled by Claude Sonnet 5.5 (claude-sonnet-5-5) for Andr\'as Kornai}
\date{}
\begin{document}
\maketitle

\section{Introduction}
North Awyu (ISO 639-3 \texttt{yir}, Glottocode nort2918) is a poorly documented Trans-New Guinea language of Papua, Indonesia. Wikipedia gives about 1,500 speakers, citing a 1987 figure. The New Guinea World database states that no North Awyu data is available to it, so its placement rests on lexical similarity figures from earlier survey work. No recent speaker-count update was found.

\section{Primary Educational Resources}
No primers, textbooks or school programmes were found.

\section{Linguistic and Reference Materials}
\begin{itemize}
\item North Awyu, Wikipedia: \url{https://en.wikipedia.org/wiki/North_Awyu}
\item Awyu languages, Wikipedia: \url{https://en.wikipedia.org/wiki/Awyu_languages}
\item Awyu people, Wikipedia: \url{https://en.wikipedia.org/wiki/Awyu_people}
\item New Guinea World, Awyu family page: \url{https://newguineaworld.linguistik.uzh.ch/families/trans-new-guinea/central-west-new-guinea/digul-river-ok/digul-river/central-digul-river/awyu/start}
\item The Awyu-Ndumut Languages (VU Amsterdam project page; it covers related Awyu languages and was not confirmed to contain North Awyu data): \url{https://www2.let.vu.nl/awyu-ndumut}
\item de Vries 2012, Glottolog reference entry: \url{https://glottolog.org/resource/reference/id/305807}
\end{itemize}

\section{Religious and Community Resources}
None found.

\section{Digital and Audio Resources}
None found.

\section{Curriculum Development Context}
There is essentially no documentation of North Awyu. Related Awyu languages have been studied (see the VU Amsterdam project above), but any curriculum would require new documentation and community work first.
\end{document}